\reading{IM-17}{Illiquid Assets}{QUICK WIN}{2}

\srcnote{Andrew Ang, \textit{Asset Management: A Systematic Approach to Factor
Investing} (New York, NY: Oxford University Press, 2014), Chapter 13. The GARP
text exists but sits inside the \emph{Liquidity} source file as its chapter 19,
because the 2026 syllabus moved this reading from LTR into IM.}

\begin{imtrapbox}[breakable=false,title={Three copies, three numbering schemes, two wrong areas}]
This reading appears \emph{three times} in the source notes. The crash layer
(labelled \emph{82}) runs page headers for \term{two different wrong exam areas}
--- ``Operational Risk Measurement \& Management'' and ``Liquidity And Treasury
Risk Measurement And Management'' --- on the same pages. The lecture layer appears
twice: once labelled \emph{99}, which is correct, and once labelled \emph{77}, a
duplicate carrying Symbol-font bullets and an internally repeated passage. The
\term{99-labelled block is canonical}; the 77-labelled duplicate is retired, and
the crash layer is merged into it.
\end{imtrapbox}

\begin{imkeybox}[title={The six objectives in this reading}]
\begin{enumerate}[label=\textbf{\alph*.}]
\item Evaluate the characteristics of illiquid markets.
\item Discuss the relationship between market imperfections and illiquidity.
\item Assess the impact of biases on reported returns for illiquid assets.
\item Explain the process of unsmoothing returns and the effects of unsmoothing.
\item Compare illiquidity risk premiums across and within asset categories.
\item Evaluate the impact of allocating illiquid assets to a portfolio, including the impact on rebalancing and trading and on optimizing the proportion of illiquid assets.
\end{enumerate}
\end{imkeybox}

% ------------------------------------------------------------------
\lo{IM-17 a}{Evaluate the characteristics of illiquid markets.}

\begin{enumerate}
\item \term{Most asset classes are illiquid.} All markets, including large-cap
equity markets, exhibit some degree of illiquidity. Illiquid assets trade
infrequently, in small amounts, with low turnover. Public equities trade rapidly;
\term{real estate and art transactions can be decades apart}.
\item \term{Illiquid asset markets are large.} The U.S.\ residential mortgage
market and the institutional real estate market each \term{surpass the market
capitalization of major stock exchanges}. Illiquid asset classes as a whole are
larger than the traditional liquid public markets of stocks and bonds.
\item \term{Investors hold lots of illiquid assets.} Homes and fine art make up a
significant portion of individual and institutional portfolios. \term{Most
investors' wealth is tied up in illiquid assets}, so asset owners must consider
illiquidity risk in portfolio construction.
\item \term{Liquidity dries up.} Even normally liquid markets periodically become
illiquid.
\end{enumerate}

% ------------------------------------------------------------------
\lo{IM-17 b}{Discuss the relationship between market imperfections and illiquidity.}

\renewcommand{\arraystretch}{1.2}
\noindent\begin{tabularx}{\textwidth}{@{}>{\raggedright\arraybackslash}p{3.2cm}X@{}}
\toprule
Clientele effects / participation cost & Entering markets can be costly; investors often must spend money, time or energy to learn their way and gain the necessary skills. \\
Transaction costs & Commissions, taxes, costs of due diligence, fees paid to lawyers, accountants and investment bankers. \\
Search frictions & For many illiquid assets you must \term{search to find an appropriate buyer or seller}, so you may wait a long time to transact. \\
Asymmetric information & One investor has superior knowledge. \term{Fearing they will be fleeced, investors become reluctant to trade.} \\
Price impact & Large trades move markets. \\
Funding constraints & Many vehicles used to invest in illiquid assets are highly levered, so a funding shock forces sales. \\
\bottomrule
\end{tabularx}
\renewcommand{\arraystretch}{1}
\figcap{Six imperfections, only one of which is a fee. Search frictions and
asymmetric information are the two that cannot be priced away --- they make the
\emph{counterparty} hard to find or hard to trust, which no amount of willingness
to pay commissions resolves.}

% ------------------------------------------------------------------
\lo{IM-17 c}{Assess the impact of biases on reported returns for illiquid assets.}

\begin{enumerate}
\item \term{Survivorship bias.} Poorly performing funds stop reporting, and many
ultimately fail --- but we rarely count their failures. In private equity and hedge
funds there is no requirement to report, so weak funds stop reporting or never
report at all. \term{Reported returns are higher than they should be.}
\item \term{Sample selection bias.} Asset values and returns are more likely to be
reported \emph{when they are high}, because assets are often sold during peak
periods. \term{This inflates perceived returns.}
\item \term{Infrequent sampling / infrequent trading.} With infrequent trading,
estimates of risk --- volatilities, correlations and betas --- computed from
reported returns are \term{too low}. Infrequent trading \term{smooths} reported
returns, making them appear less volatile.
\end{enumerate}

\begin{imtrapbox}[breakable=false,title={Two biases inflate returns, one deflates risk}]
Survivorship and sample selection both push \term{reported returns up}. Infrequent
trading pushes \term{reported risk down}. All three flatter the asset class, but
they act on different halves of the risk-return pair --- and only the third is
fixable by a filter.
\end{imtrapbox}

% ------------------------------------------------------------------
\lo{IM-17 d}{Explain the process of unsmoothing returns and the effects of unsmoothing.}

\begin{imdefbox}[title={Unsmoothing}]
To account for the infrequent trading bias we need to go from samples that are too
smooth to noisier data. This process is called \term{unsmoothing} or
\term{de-smoothing}, and it is a \term{filter problem}.
\end{imdefbox}

\begin{imkeybox}[title={Two properties worth memorising}]
\begin{enumerate}
\item \term{Unsmoothing affects only risk estimates, not expected returns.}
Intuitively, estimates of the mean require only the first and last price
observation, which unsmoothing leaves alone.
\item \term{Unsmoothing has no effect if the observed returns are uncorrelated.}
It works by undoing autocorrelation; with none present, there is nothing to undo.
\end{enumerate}
\end{imkeybox}

\begin{imtrapbox}[breakable=false,title={The most corruptible pair of facts in the reading}]
``Unsmoothing raises the estimated return'' is false --- it raises the estimated
\emph{risk} and leaves the mean where it was. And ``unsmoothing always increases
volatility'' is false too: with uncorrelated observed returns it does nothing at
all.
\end{imtrapbox}

% ------------------------------------------------------------------
\lo{IM-17 e}{Compare illiquidity risk premiums across and within asset categories.}

\subsection{Across asset classes}

There is a positive connection between the illiquidity of an asset class and its
anticipated return. \term{Venture capital}, the least liquid, has the highest
expected return at around \term{16--17\%}, while more liquid assets such as
equities and cash offer roughly \term{4\%}.

\begin{imtrapbox}[breakable=false,title={The conventional view is challenged}]
Three reasons the across-class premium may not be real:
\begin{enumerate}
\item \term{Illiquidity biases.} Illiquid assets have \emph{overstated} returns
using raw data and \emph{underestimated} risks and correlations --- exactly the
three biases of IM-17 c.
\item \term{Additional risks.} Private equity and timber carry risks beyond
liquidity which, once considered, make them less attractive.
\item \term{No market index.} Unlike liquid assets, there is no investable index,
so the reported ``asset class return'' is not something an investor could have
earned.
\end{enumerate}
\end{imtrapbox}

\subsection{Within asset classes}

Less liquid assets tend to offer \term{higher returns than more liquid assets
within the same asset class}. There is \term{no established theory} explaining why
the premium exists \emph{within} classes but not \emph{between} them. Possible
reasons include investor overvaluation of illiquid asset classes and institutional
factors that limit arbitrage.

\begin{imkeybox}[title={Illiquidity effects in other markets}]
\term{U.S.\ Treasuries.} Newly issued \term{on-the-run} bills are more liquid and
have \emph{lower} yields than seasoned \term{off-the-run} bills --- the
\term{on-the-run/off-the-run spread}. During 2007--2009, same-maturity Treasury
bonds and notes traded at different yields, with \term{T-bond prices more than 5\%
lower} than T-note prices.

\smallskip
\term{Private equity secondaries.} PE funds trade portfolio companies with each
other, about \term{15\% of buyout deals in 2005}. This allows exit from specific
deals but \term{not from the fund itself}. LPs exiting funds face immature, small
and less transparent secondary markets. In the 1990s buy-side firms were called
\term{``vultures''}, securing discounts of \term{30\% to 50\%}; discounts fell
below \term{20\%} in the early 2000s before widening again in the crisis.
\end{imkeybox}

\subsection{Harvesting the premium}

\begin{enumerate}
\item \term{Passive allocation} to illiquid asset classes such as real estate.
\item \term{Liquidity security selection} --- choosing more illiquid assets within
an asset class.
\item Acting as a \term{market maker} for individual securities.
\item \term{Dynamic factor strategies} at the aggregate portfolio level: long
illiquid assets, short liquid assets.
\end{enumerate}

\begin{imkeybox}
Of the four, the \term{fourth is the easiest to implement and can have the
greatest effect on portfolio returns} --- which is the ranking the reading asks
you to remember, and it is not the one intuition suggests.
\end{imkeybox}

% ------------------------------------------------------------------
\lo{IM-17 f}{Evaluate the impact of allocating illiquid assets to a portfolio, including the impact on rebalancing and trading and on optimizing the proportion of illiquid assets.}

\srcnote{The source notes head this objective ``Evaluate portfolio choice
decisions on the inclusion of illiquid assets,'' which drops both of the spine's
named halves --- the impact on rebalancing and trading, and the optimization of
the proportion held. The spine wording is restored above.}

Portfolio choice models that include illiquid assets must consider two aspects of
illiquidity that impact investors: \term{long time horizons between trades} and
\term{large transaction costs}.

\begin{imkeybox}[title={Why standard asset allocation breaks}]
Asset allocation models often assume assets can \emph{always} be traded if
transaction costs are covered. That assumption does not hold in private equity,
infrastructure, real estate and timber, where \term{finding a buyer can be
challenging and time-consuming}. During market stress even liquid asset classes
may face liquidity problems.
\end{imkeybox}

\subsection{Consequences for the investor}

\begin{enumerate}
\item \term{Infrequent trading.} Investors cannot easily adjust holdings, making
\term{portfolio rebalancing difficult}.
\item \term{Lower optimal holdings.} The optimal allocation to illiquid assets
should be \term{lower} than for liquid ones.
\item \term{No liquidity hedge.} Investors cannot hedge against declining value
because they cannot easily sell.
\item \term{Illiquidity risk premium.} Investors need compensation for these
risks, typically through higher expected returns.
\end{enumerate}

\subsection{Six cautions on including illiquid assets}

\begin{enumerate}
\item Studies show illiquid assets \term{do not deliver higher risk-adjusted
returns}.
\item Investors are subject to \term{agency problems}, relying on the talents of a
manager who is difficult to monitor.
\item In many firms, illiquid assets are \term{managed separately} from the rest
of the portfolio.
\item Illiquid asset markets are \term{less efficient}, with information
asymmetries and high transaction costs deterring many investors.
\item Success requires \term{specialized skills and resources}, and the ability to
find, evaluate and monitor opportunities.
\item Endowments such as \term{Harvard, Yale and Stanford} have the expertise and
resources to navigate these markets effectively --- which is the point: most
investors do not.
\end{enumerate}

\begin{imtrapbox}[title={Consolidated trap summary --- IM-17}]
\begin{itemize}
\item \term{Definition, IM-17 a.} Illiquid markets are \emph{large} --- larger in
aggregate than public stock and bond markets --- not niche.
\item \term{Sibling, IM-17 c.} Survivorship and selection bias inflate
\emph{returns}; infrequent trading deflates \emph{risk}. Different halves of the
pair.
\item \term{Polarity, IM-17 d.} Unsmoothing changes \emph{risk} estimates only,
never expected returns, and has \emph{no effect at all} if observed returns are
uncorrelated.
\item \term{Number, IM-17 e.} Venture capital around 16--17\%; equities and cash
around 4\%.
\item \term{Polarity, IM-17 e.} On-the-run Treasuries are \emph{more} liquid and
carry \emph{lower} yields. Reversing this reverses the whole illiquidity argument.
\item \term{Scope, IM-17 e.} PE secondaries let you exit a \emph{deal}, not the
\emph{fund}. LP exits face a separate, immature market.
\item \term{Ranking, IM-17 e.} The \emph{dynamic factor strategy} is the easiest
of the four harvesting methods to implement and has the greatest potential effect.
\item \term{Polarity, IM-17 f.} Optimal allocation to illiquid assets is
\emph{lower} than to liquid ones, and studies find \emph{no} higher risk-adjusted
return.
\end{itemize}
\end{imtrapbox}
